\chapter{Curvature}
\label{chap:pet-curvature}

Throughout, $\nabla$ denotes the Levi--Civita connection of $(M,g)$. Curvature
conventions are fixed by \cref{def:pet-ch3-curvature-tensor}.

\section{Curvature}

The curvature tensor is first viewed as a commutator of covariant derivatives,
then contracted or transferred to $\Lambda^2TM$ to obtain its standard scalar
invariants.

\subsection{The Curvature Tensor}

The sign convention is fixed by the following definition.

\begin{definition}[curvature tensor]
  \label{def:pet-ch3-curvature-tensor}
  \lean{PetersenLib.curvatureTensor}
  \source{petersen:page-0096}
  \leanok
  Let $(M,g)$ be a Riemannian manifold with Riemannian connection $\nabla$. The
  \emph{curvature tensor} is the $(1,3)$-tensor defined on vector fields $X,Y,Z$ by
  \[
    R(X,Y)Z=\nabla^2_{X,Y}Z-\nabla^2_{Y,X}Z
    =\nabla_X\nabla_YZ-\nabla_Y\nabla_XZ-\nabla_{[X,Y]}Z
    =[\nabla_X,\nabla_Y]Z-\nabla_{[X,Y]}Z.
  \]
  The first expression is also called the \emph{Ricci identity} and is written
  $R_{X,Y}Z=\nabla^2_{X,Y}Z-\nabla^2_{Y,X}Z$. The same formula defines the
  curvature of an arbitrary tensor $T$:
  \[
    R(X,Y)T=R_{X,Y}T=\nabla^2_{X,Y}T-\nabla^2_{Y,X}T.
  \]
\end{definition}

\begin{lemma}[tensoriality of the curvature tensor]
  \label{lem:pet-ch3-curvature-tensor-tensorial}
  \lean{PetersenLib.curvatureTensor_tensorial}
  \source{petersen:page-0096}
  \uses{def:pet-ch3-curvature-tensor}
  \leanok
  For $f\in C^\infty(M)$,
  $R(X,Y)(fZ)=fR(X,Y)Z$.
\end{lemma}
\begin{proof}
  \uses{lem:pet-ch3-curvature-tensor-tensorial}
  \leanok
  \sketch
  The product rule gives
  \[
   R(X,Y)(fZ)=fR(X,Y)Z+
   \bigl(\nabla^2_{X,Y}f-\nabla^2_{Y,X}f\bigr)Z.
  \]
  The final term vanishes because the Hessian of a function is symmetric.
\end{proof}

\begin{remark}[the Ricci identity as a derivation]
  \label{rem:pet-ch3-ricci-identity-derivation}
  \lean{PetersenLib.curvatureTensor_derivation}
  \source{petersen:page-0097}
  \uses{lem:pet-ch3-curvature-tensor-tensorial}
  \leanok
  The symmetry of the Hessian implies that $R_{X,Y}$ annihilates functions. Thus
  $\nabla^2_{X,Y}-\nabla^2_{Y,X}=R_{X,Y}=R(X,Y)$, where on the right $R(X,Y)$ is
  regarded as a $(1,1)$-tensor acting on tensors as a derivation. In particular, if
  $T$ is a $(0,k)$-tensor,
  \[
    (R_{X,Y}T)(X_1,\dots,X_k)=-T(R(X,Y)X_1,\dots,X_k)-\dots-T(X_1,\dots,R(X,Y)X_k).
  \]
\end{remark}

\begin{definition}[$(0,4)$-curvature tensor]
  \label{def:pet-ch3-curvature-04-tensor}
  \lean{PetersenLib.curvatureTensorFour}
  \source{petersen:page-0097}
  \uses{def:pet-ch3-curvature-tensor}
  \leanok
  Lower the output index with $g$:
  \[
    R(X,Y,Z,W)=g(R(X,Y)Z,W).
  \]
\end{definition}

\begin{proposition}[Symmetries of the curvature tensor]
  \label{prop:pet-ch3-curvature-symmetries}
  \lean{PetersenLib.curvatureTensor_symmetries}
  \source{petersen:page-0097,petersen:page-0098,petersen:page-0099,petersen:page-0100}
  \uses{def:pet-ch3-curvature-04-tensor}
  \leanok
  \dcref{ch3:3.1.1}
  The Riemannian curvature tensor $R(X,Y,Z,W)$ satisfies:
  \begin{enumerate}
    \item[(1)] skew-symmetry in the first two and in the last two entries:
    $R(X,Y,Z,W)=-R(Y,X,Z,W)=R(Y,X,W,Z)$;
    \item[(2)] symmetry between the first two and the last two entries:
    $R(X,Y,Z,W)=R(Z,W,X,Y)$;
    \item[(3)] Bianchi's first identity: $R(X,Y)Z+R(Z,X)Y+R(Y,Z)X=0$;
    \item[(4)] Bianchi's second identity:
    $(\nabla_ZR)_{X,Y}W+(\nabla_XR)_{Y,Z}W+(\nabla_YR)_{Z,X}W=0$.
  \end{enumerate}
\end{proposition}
\begin{proof}
  \uses{prop:pet-ch3-curvature-symmetries}
  \leanok
  \sketch
  Skew-symmetry in $X,Y$ follows from the definition. Applying
  $R_{X,Y}$ to $g(Z,Z)$ and using metric compatibility yields
  $R(X,Y,Z,Z)=0$; polarization gives skew-symmetry in $Z,W$.
  Torsion-freeness and the Jacobi identity for Lie brackets give the first
  Bianchi identity, which together with the two skew-symmetries implies pair
  symmetry. Finally, expand the cyclic sum of $\nabla R$: third derivatives and
  connection-correction terms cancel in pairs, while the remaining bracket terms
  vanish by the Jacobi identity. This is the second Bianchi identity.
\end{proof}

\begin{example}[Euclidean space is flat]
  \label{ex:pet-ch3-euclidean-flat}
  \lean{PetersenLib.euclideanSpace_curvature_eq_zero}
  \source{petersen:page-0100}
  \uses{def:pet-ch3-curvature-tensor}
  \leanok
  \dcref{ch3:3.1.2}
  $(\mathbb R^n,g_{\mathrm{eu}})$ has $R\equiv0$, since $\nabla_i\partial_j=0$ for the
  standard Cartesian coordinate vector fields.
\end{example}

\subsection{The Curvature Operator}

The pair symmetries of $R$ allow curvature to be encoded by a self-adjoint
operator on bivectors.

\begin{definition}[inner product on bivectors]
  \label{def:pet-ch3-bivector-inner-product}
  \lean{PetersenLib.bivectorInnerProduct}
  \source{petersen:page-0100}
  \leanok
  On the space $\Lambda^2T_pM$ of bivectors, with $e_1,\dots,e_n$ an orthonormal
  basis of $T_pM$, the inner product making $\{e_i\wedge e_j\}_{i<j}$ an
  orthonormal basis is denoted $g$ as well, and satisfies
  \[
    g(x\wedge y,v\wedge w)=g(x,v)g(y,w)-g(x,w)g(y,v)
    =\det\begin{pmatrix}g(x,v)&g(x,w)\\g(y,v)&g(y,w)\end{pmatrix}.
  \]
\end{definition}

\begin{definition}[bivectors as skew-symmetric maps]
  \label{def:pet-ch3-bivector-skew-map}
  \lean{PetersenLib.bivectorSkewMap}
  \source{petersen:page-0100}
  \uses{def:pet-ch3-bivector-inner-product}
  \leanok
  Identify a decomposable bivector with the skew-adjoint map
  \[
    (x\wedge y)(v)=g(x,v)y-g(y,v)x,
  \]
  which satisfies
  $g(x\wedge y,v\wedge w)=g(x,v)g(y,w)-g(x,w)g(y,v)=g((x\wedge y)(v),w)$.
\end{definition}

\begin{lemma}[Jacobi identity for bivector maps]
  \label{lem:pet-ch3-bivector-jacobi-identity}
  \lean{PetersenLib.bivectorSkewMap_jacobi}
  \source{petersen:page-0100}
  \uses{def:pet-ch3-bivector-skew-map}
  \leanok
  $(x\wedge y)(z)+(y\wedge z)(x)+(z\wedge x)(y)=0$.
\end{lemma}
\begin{proof}
  \uses{lem:pet-ch3-bivector-jacobi-identity}
  \leanok
  \sketch
  Expanding the three skew maps gives
  \[
   \bigl(g(x,z)-g(z,x)\bigr)y+
   \bigl(g(y,x)-g(x,y)\bigr)z+
   \bigl(g(z,y)-g(y,z)\bigr)x=0.
  \]
\end{proof}

\begin{definition}[curvature operator]
  \label{def:pet-ch3-curvature-operator}
  \lean{PetersenLib.curvatureOperator}
  \source{petersen:page-0101}
  \uses{def:pet-ch3-curvature-04-tensor, def:pet-ch3-bivector-inner-product, def:pet-ch3-bivector-skew-map}
  \leanok
  By the symmetry properties of $R$, it defines a symmetric bilinear map
  $R:\Lambda^2TM\times\Lambda^2TM\to\mathbb R$,
  \[
    R\Big(\sum X_i\wedge Y_i,\sum V_j\wedge W_j\Big)=\sum R(X_i,Y_i,W_j,V_j)
  \]
  (note the reversal of $V_j,W_j$). Consequently there is a unique self-adjoint
  operator $\mathfrak R:\Lambda^2TM\to\Lambda^2TM$, the \emph{curvature operator},
  with
  \[
    g\Big(\mathfrak R\Big(\sum X_i\wedge Y_i\Big),\sum V_j\wedge W_j\Big)
    =\sum R(X_i,Y_i,W_j,V_j).
  \]
\end{definition}

\subsection{Sectional Curvature}

Restricting the curvature operator to decomposable bivectors defines curvature
of tangent $2$-planes.

\begin{definition}[directional curvature operator]
  \label{def:pet-ch3-directional-curvature-operator}
  \lean{PetersenLib.directionalCurvatureOperator}
  \source{petersen:page-0101}
  \uses{def:pet-ch3-curvature-tensor, prop:pet-ch3-curvature-symmetries}
  \leanok
  For $v\in T_pM$, the \emph{directional} (or \emph{tidal force}) \emph{curvature
  operator} is $R_v(w)=R(w,v)v:T_pM\to T_pM$. It is self-adjoint and $v$ is always
  a zero eigenvector, by the symmetry properties of $R$.
\end{definition}

\begin{definition}[sectional curvature]
  \label{def:pet-ch3-sectional-curvature}
  \lean{PetersenLib.sectionalCurvature}
  \source{petersen:page-0101}
  \uses{def:pet-ch3-directional-curvature-operator, def:pet-ch3-bivector-inner-product}
  \leanok
  The \emph{sectional curvature} of $v,w\in T_pM$ is the normalized biquadratic form
  \[
    \sec(v,w)=\frac{g(R_v(w),w)}{g(v,v)g(w,w)-g(v,w)^2}
    =\frac{g(R(w,v)v,w)}{g(v\wedge w,v\wedge w)}.
  \]
  Since the denominator is the squared area of the parallelogram spanned by
  $v,w$, $\sec(v,w)$ depends only on the plane $\pi=\operatorname{span}\{v,w\}$.
\end{definition}

\begin{proposition}[Riemann's characterization of constant sectional curvature]
  \label{prop:pet-ch3-riemann-constant-curvature-equivalence}
  \lean{PetersenLib.riemann_constantCurvature_equivalence}
  \source{petersen:page-0101,petersen:page-0102,petersen:page-0103}
  \uses{def:pet-ch3-sectional-curvature, def:pet-ch3-curvature-operator, prop:pet-ch3-curvature-symmetries, lem:pet-ch3-bivector-jacobi-identity}
  \leanok
  \dcref{ch3:3.1.3}
  The following are equivalent, for fixed $p\in M$ and $k\in\mathbb R$:
  \begin{enumerate}
    \item[(1)] $\sec(\pi)=k$ for all $2$-planes $\pi\subset T_pM$;
    \item[(2)] $R(v_1,v_2)v_3=-k(v_1\wedge v_2)(v_3)$ for all $v_1,v_2,v_3\in T_pM$;
    \item[(3)] $R_v(w)=k\cdot(w-g(w,v)v)=k\cdot\operatorname{prj}_{v^\perp}(w)$ for all
    $w\in T_pM$ and unit $v$;
    \item[(4)] $\mathfrak R(\omega)=k\cdot\omega$ for all $\omega\in\Lambda^2T_pM$.
  \end{enumerate}
\end{proposition}
\begin{proof}
  \uses{prop:pet-ch3-riemann-constant-curvature-equivalence}
  \leanok
  \sketch
  By \cref{def:pet-ch3-directional-curvature-operator} and
  \cref{def:pet-ch3-sectional-curvature},
  conditions (2) and (3) are equivalent by evaluating the bivector map on a
  unit vector, and either gives (1) after pairing with $w$. Conversely, compare
  $R$ with
  \[
   R_k(x,y)z=-k(x\wedge y)(z).
  \]
  Their difference has the curvature symmetries; the first Bianchi identity for
  $R_k$ follows from \cref{lem:pet-ch3-bivector-jacobi-identity}. It vanishes on
  every expression $D(v,w,w,v)$, so polarization and Bianchi force $D=0$,
  proving (2). The defining pairing of $\mathfrak R$ in
  \cref{def:pet-ch3-curvature-operator} then makes (2) equivalent to
  $\mathfrak R=k\,\mathrm{id}$, which is (4).
\end{proof}

\begin{definition}[constant curvature]
  \label{def:pet-ch3-constant-curvature}
  \lean{PetersenLib.HasConstantCurvature}
  \source{petersen:page-0103}
  \uses{prop:pet-ch3-riemann-constant-curvature-equivalence}
  \leanok
  $(M,g)$ has \emph{constant curvature} $k\in\mathbb R$ if one (equivalently, by
  \cref{prop:pet-ch3-riemann-constant-curvature-equivalence}, all) of the four
  equivalent conditions holds at every $p\in M$ with the same $k$.
\end{definition}

\subsection{Ricci Curvature}

Tracing one input-output pair of $R$ produces Ricci curvature.

\begin{definition}[Ricci curvature]
  \label{def:pet-ch3-ricci-curvature}
  \lean{PetersenLib.RicciCurvature}
  \source{petersen:page-0103}
  \uses{def:pet-ch3-curvature-tensor}
  \leanok
  For an orthonormal basis $e_1,\dots,e_n$ of $T_pM$, the \emph{Ricci curvature} is
  the trace
  \[
    \operatorname{Ric}(v,w)=\operatorname{tr}(x\mapsto R(x,v)w)
    =\sum_{i=1}^ng(R(e_i,v)w,e_i)=\sum_{i=1}^ng(R(v,e_i)e_i,w)
    =\sum_{i=1}^ng(R(e_i,w)v,e_i),
  \]
  a symmetric bilinear form, equivalently the symmetric $(1,1)$-tensor
  $\operatorname{Ric}(v)=\sum_{i=1}^nR(v,e_i)e_i$.
\end{definition}

\begin{definition}[Einstein manifold]
  \label{def:pet-ch3-einstein-manifold}
  \lean{PetersenLib.IsEinstein}
  \source{petersen:page-0103,petersen:page-0104}
  \uses{def:pet-ch3-ricci-curvature}
  \leanok
  Write $\operatorname{Ric}\ge k$ if, at every $p\in M$, each eigenvalue of the
  Ricci endomorphism $\operatorname{Ric}_p:T_pM\to T_pM$ is $\ge k$; equivalently,
  $\operatorname{Ric}(v,v)\ge kg(v,v)$ for every tangent vector $v$. $(M,g)$ is an
  \emph{Einstein manifold with Einstein constant $k$} if $\operatorname{Ric}(v)=k\cdot v$
  for all $v$, equivalently $\operatorname{Ric}(v,w)=k\cdot g(v,w)$.
\end{definition}

\begin{proposition}[Ricci curvature via sectional curvature]
  \label{prop:pet-ch3-ricci-sectional-relation}
  \lean{PetersenLib.ricciCurvature_eq_sum_sectionalCurvature}
  \source{petersen:page-0104}
  \uses{def:pet-ch3-ricci-curvature, def:pet-ch3-sectional-curvature}
  \leanok
  If $v\in T_pM$ is a unit vector completed to an orthonormal basis
  $\{v,e_2,\dots,e_n\}$, then
  $\operatorname{Ric}(v,v)=\sum_{i=2}^n\sec(v,e_i)$.
\end{proposition}
\begin{proof}
  \uses{prop:pet-ch3-ricci-sectional-relation}
  \leanok
  \sketch
  By \cref{def:pet-ch3-ricci-curvature}, complete $v$ to the stated orthonormal
  basis. In the trace defining
  $\operatorname{Ric}(v,v)$, the $v$-term is zero and each remaining term is
  $g(R(e_i,v)v,e_i)=\sec(v,e_i)$ by
  \cref{def:pet-ch3-sectional-curvature}.
\end{proof}

\begin{corollary}[Ricci determines sec in dimensions $2$ and $3$]
  \label{cor:pet-ch3-ricci-determines-sec-low-dim}
  \lean{PetersenLib.einstein_iff_hasConstantCurvature_of_finrank_three}
  \source{petersen:page-0104}
  \uses{prop:pet-ch3-ricci-sectional-relation, def:pet-ch3-einstein-manifold, def:pet-ch3-constant-curvature}
  \leanok
  For $n=2$, and also for $n=3$, $\operatorname{Ric}$ determines $\sec$
  completely: for $n=3$ and orthonormal $\{e_1,e_2,e_3\}$,
  \[
    \begin{pmatrix}1&0&1\\1&1&0\\0&1&1\end{pmatrix}
    \begin{pmatrix}\sec(e_1,e_2)\\\sec(e_2,e_3)\\\sec(e_1,e_3)\end{pmatrix}
    =\begin{pmatrix}\operatorname{Ric}(e_1,e_1)\\\operatorname{Ric}(e_2,e_2)\\\operatorname{Ric}(e_3,e_3)\end{pmatrix},
  \]
  a system with determinant $2$; hence a $3$-manifold is Einstein if and only
  if it has constant sectional curvature.
\end{corollary}
\begin{proof}
  \uses{cor:pet-ch3-ricci-determines-sec-low-dim}
  \leanok
  \sketch
  In dimension $3$, the three diagonal Ricci values are the pairwise sums of
  the three sectional curvatures displayed in the statement. The coefficient
  matrix has determinant $2$, so these curvatures are uniquely recovered. The
  dimension-$2$ assertion is immediate from the single tangent $2$-plane.
\end{proof}

\begin{remark}[constant curvature is Einstein]
  \label{rem:pet-ch3-constant-curvature-einstein}
  \lean{PetersenLib.constantCurvature_isEinstein}
  \source{petersen:page-0104}
  \uses{def:pet-ch3-constant-curvature, def:pet-ch3-einstein-manifold, prop:pet-ch3-ricci-sectional-relation}
  \leanok
  \sketch
  If $(M,g)$ has constant curvature $k$, it is Einstein with Einstein constant
  $(n-1)k$: for unit $v$, \cref{prop:pet-ch3-ricci-sectional-relation} gives
  $\operatorname{Ric}(v,v)=\sum_{i=2}^n\sec(v,e_i)=(n-1)k$.
\end{remark}

\begin{remark}[dimension of nonconstant-curvature Einstein metrics]
  \label{rem:pet-ch3-einstein-not-constant-curvature-examples}
  \lean{PetersenLib.einstein_examples}
  \source{petersen:page-0104}
  \uses{def:pet-ch3-einstein-manifold, prop:pet-ch3-ricci-sectional-relation}
  % Lean: only the dimension-$3$ obstruction is formalized; the metric families are omitted.
  \sketch
  An Einstein metric without constant sectional curvature must have dimension
  at least $4$, by
  \cref{cor:pet-ch3-ricci-determines-sec-low-dim}.
\end{remark}

\subsection{Scalar Curvature}

The full trace of curvature is the scalar curvature.

\begin{definition}[scalar curvature]
  \label{def:pet-ch3-scalar-curvature}
  \lean{PetersenLib.scalarCurvature}
  \source{petersen:page-0104,petersen:page-0105}
  \uses{def:pet-ch3-ricci-curvature, def:pet-ch3-curvature-operator}
  \leanok
  The \emph{scalar curvature} is $\operatorname{scal}=\operatorname{tr}(\operatorname{Ric})=2\operatorname{tr}\mathfrak R$,
  a function $\operatorname{scal}:M\to\mathbb R$.
\end{definition}

\begin{remark}[formulas for the scalar curvature]
  \label{rem:pet-ch3-scalar-curvature-formulas}
  \lean{PetersenLib.scalarCurvature_formulas}
  \source{petersen:page-0105}
  \uses{def:pet-ch3-scalar-curvature}
  \leanok
  In an orthonormal basis $e_1,\dots,e_n$,
  \[
    \operatorname{scal}=\sum_{j}g(\operatorname{Ric}(e_j),e_j)
    =\sum_{i,j}g(R(e_i,e_j)e_j,e_i)=\sum_{i,j}g(\mathfrak R(e_i\wedge e_j),e_i\wedge e_j)
    =2\sum_{i<j}g(\mathfrak R(e_i\wedge e_j),e_i\wedge e_j)=2\sum_{i<j}\sec(e_i,e_j).
  \]
  When $n=2$, $\operatorname{scal}(p)=2\sec(T_pM)$.
\end{remark}

\begin{notation}[divergence/adjoint of a tensor]
  \label{notation:pet-ch3-divergence-adjoint}
  \lean{PetersenLib.divergenceAdjoint}
  \source{petersen:page-0105,petersen:page-0106}
  \leanok
  For a $(0,k)$-tensor $T$ and an orthonormal frame $E_1,\dots,E_n$ on $(M,g)$, the
  adjoint of the covariant derivative is the $(0,k-1)$-tensor
  \[
    (\nabla^*T)(X_1,\dots,X_{k-1})=-\sum_{i=1}^n(\nabla_{E_i}T)(E_i,X_1,\dots,X_{k-1}),
  \]
  independent of the choice of orthonormal frame.
\end{notation}

\begin{proposition}[Contracted Bianchi identity]
  \label{prop:pet-ch3-contracted-bianchi}
  \lean{PetersenLib.contractedBianchiIdentity}
  \source{petersen:page-0106}
  \uses{def:pet-ch3-scalar-curvature, def:pet-ch3-ricci-curvature, notation:pet-ch3-divergence-adjoint, prop:pet-ch3-curvature-symmetries}
  \leanok
  \dcref{ch3:3.1.5}
  On any Riemannian manifold, $d\operatorname{tr}(\operatorname{Ric})=d\operatorname{scal}=-2\nabla^*\operatorname{Ric}$.
\end{proposition}
\begin{proof}
  \uses{prop:pet-ch3-contracted-bianchi}
  \leanok
  \sketch
  At a fixed point choose a normal orthonormal frame $(E_i)$. Differentiate
  $\operatorname{scal}=\sum_{i,j}R(E_i,E_j,E_j,E_i)$, apply the second Bianchi
  identity, and use the curvature symmetries to identify the two resulting
  contractions using \cref{def:pet-ch3-ricci-curvature}. This gives
  \[
   d\operatorname{scal}(W)
   =2\sum_i(\nabla_{E_i}\operatorname{Ric})(E_i,W)
   =-2(\nabla^*\operatorname{Ric})(W),
  \]
  where the last equality is the convention of
  \cref{notation:pet-ch3-divergence-adjoint}.
\end{proof}

\begin{lemma}[Schur's lemma]
  \label{lem:pet-ch3-schur-lemma}
  \lean{PetersenLib.schurLemma}
  \source{petersen:page-0105,petersen:page-0106}
  \uses{def:pet-ch3-sectional-curvature, def:pet-ch3-einstein-manifold, prop:pet-ch3-contracted-bianchi, prop:pet-ch3-riemann-constant-curvature-equivalence}
  \leanok
  \dcref{ch3:3.1.4}
  Let $(M,g)$ have dimension $n\ge3$ and satisfy, for a function $f:M\to\mathbb
  R$, either
  \begin{enumerate}
    \item[(1)] $\sec(\pi)=f(p)$ for all $2$-planes $\pi\subset T_pM$, $p\in M$, or
    \item[(2)] $\operatorname{Ric}(v)=(n-1)f(p)v$ for all $v\in T_pM$, $p\in M$.
  \end{enumerate}
  Then $f$ is constant; the metric has constant curvature, respectively is
  Einstein.
\end{lemma}
\begin{proof}
  \uses{lem:pet-ch3-schur-lemma}
  \leanok
  \sketch
  The sectional-curvature hypothesis implies
  $\operatorname{Ric}=(n-1)fg$, so it suffices to treat the Ricci hypothesis.
  Taking traces gives $\operatorname{scal}=n(n-1)f$, whereas metric compatibility
  gives $\nabla^*\operatorname{Ric}=-(n-1)df$.
  The identity \cref{prop:pet-ch3-contracted-bianchi} therefore yields
  \[
   n(n-1)df=2(n-1)df.
  \]
  Since $n\ge3$, one has $df=0$.
\end{proof}

\begin{corollary}[Einstein metrics and normalized Ricci curvature]
  \label{cor:pet-ch3-einstein-ricci-scalar-characterization}
  \lean{PetersenLib.einstein_iff_ricci_eq_scal_div_n_metric}
  \source{petersen:page-0106}
  \uses{def:pet-ch3-einstein-manifold, def:pet-ch3-scalar-curvature, lem:pet-ch3-schur-lemma}
  \leanok
  \dcref{ch3:3.1.6}
  An $n(>2)$-dimensional Riemannian manifold $(M,g)$ is Einstein if and only if
  $\operatorname{Ric}=\dfrac{\operatorname{scal}}{n}g$.
\end{corollary}
\begin{proof}
  \uses{cor:pet-ch3-einstein-ricci-scalar-characterization}
  \leanok
  \sketch
  If $\operatorname{Ric}=kg$, then tracing gives
  $\operatorname{scal}=nk$. Conversely,
  $\operatorname{Ric}=(\operatorname{scal}/n)g$ satisfies Schur's hypothesis with
  $f=\operatorname{scal}/(n(n-1))$; \cref{lem:pet-ch3-schur-lemma} makes
  $\operatorname{scal}$, and therefore the Einstein constant, constant.
\end{proof}

\subsection{Curvature in Local Coordinates}

In a coordinate frame, curvature is expressed by first derivatives and
quadratic combinations of Christoffel symbols.

\begin{proposition}[curvature tensor in coordinates]
  \label{prop:pet-ch3-curvature-coordinate-formula}
  \lean{PetersenLib.curvatureTensor_coordinates}
  \source{petersen:page-0107}
  \uses{def:pet-ch3-curvature-tensor}
  \leanok
  In local coordinates, writing $X=X^i\partial_i$, $Y=Y^j\partial_j$,
  $Z=Z^k\partial_k$ and $R(\partial_i,\partial_j)\partial_k=R^l_{ijk}\partial_l$, so
  $R(X,Y)Z=X^iY^jZ^kR^l_{ijk}\partial_l$, the coefficients are
  \[
    R^l_{ijk}=\partial_i\Gamma^l_{jk}-\partial_j\Gamma^l_{ik}
    +\Gamma^s_{jk}\Gamma^l_{is}-\Gamma^s_{ik}\Gamma^l_{js}.
  \]
  At a point $p$ where $\Gamma^k_{ij}|_p=0$ this simplifies to
  $R^l_{ijk}|_p=\partial_i\Gamma^l_{jk}|_p-\partial_j\Gamma^l_{ik}|_p$. Using the
  formula for the Christoffel symbols in terms of $g$, $R^l_{ijk}$ depends on the
  metric coefficients $g_{ij}$ and their first two derivatives.
\end{proposition}
\begin{proof}
  \uses{prop:pet-ch3-curvature-coordinate-formula}
  \leanok
  \sketch
  Insert
  $\nabla_{\partial_i}\partial_j=\Gamma^s_{ij}\partial_s$ into
  $R(\partial_i,\partial_j)\partial_k$. The two differentiated Christoffel terms
  and the two quadratic terms are exactly the displayed formula. At a point
  where all $\Gamma^k_{ij}$ vanish, only the derivative terms remain.
\end{proof}

\begin{remark}[Ricci and scalar curvature in index notation]
  \label{rem:pet-ch3-ricci-scalar-index-notation}
  \lean{PetersenLib.ricciScalar_indexNotation}
  \source{petersen:page-0108}
  \uses{def:pet-ch3-ricci-curvature, def:pet-ch3-scalar-curvature, prop:pet-ch3-curvature-coordinate-formula}
  \leanok
  \dcref{ch3:3.1.7}
  In coordinates one writes $\operatorname{Ric}_{ij}=R_{ij}=R^k_{ik}=g^{kl}R_{kijl}$ and
  $\operatorname{scal}=R=g^{ij}R_{ij}$, both obtained by contracting indices of the curvature
  tensor; to avoid confusing the full curvature tensor with the scalar curvature
  it is also denoted $\mathfrak{Rm}$.
\end{remark}

\begin{remark}[Lowered curvature components]
  \label{rem:pet-ch3-curvature-index-conventions}
  \lean{PetersenLib.curvatureTensor_indexLowering}
  \source{petersen:page-0108}
  \uses{def:pet-ch3-curvature-04-tensor, prop:pet-ch3-curvature-coordinate-formula}
  \leanok
  \dcref{ch3:3.1.8}
  Consistently with $R(X,Y,Z,W)=g(R(X,Y)Z,W)$, the fully lower-index curvature
  components are $R_{ijkl}=g_{sl}R^s_{ijk}=R(\partial_i,\partial_j,\partial_k,\partial_l)$.
\end{remark}

\section{The Equations of Riemannian Geometry}

For regular level hypersurfaces, the ambient curvature splits into radial,
tangential, and normal components controlled by the Hessian and second
fundamental form.

\subsection{Curvature Equations}

\begin{definition}[the tensors $S$ and $\operatorname{Hess}^2f$]
  \label{def:pet-ch3-hessian-square-notation}
  \lean{PetersenLib.hessianOperator}
  \source{petersen:page-0108}
  \leanok
  For a smooth function $f$ (possibly only on an open $O\subset M$), write
  $S(X)=\nabla_X\nabla f$ for the $(1,1)$-tensor corresponding to
  $\operatorname{Hess}f$, and $\operatorname{Hess}^2f$ for the $(0,2)$-tensor
  corresponding to $S^2=S\circ S$.
\end{definition}

\begin{definition}[second fundamental form]
  \label{def:pet-ch3-second-fundamental-form}
  \lean{PetersenLib.secondFundamentalForm}
  \source{petersen:page-0108,petersen:page-0109}
  \leanok
  For a hypersurface $H^{n-1}\subset M^n$ with fixed unit normal field
  $N:H\to T^\perp H=\{v\in T_pM\mid p\in H,\,v\perp T_pH\}$, the \emph{second
  fundamental form} is the $(0,2)$-tensor $\Pi(X,Y)=g(\nabla_XN,Y)$ on $H$.
\end{definition}
\begin{proof}
  \uses{def:pet-ch3-second-fundamental-form}
  \leanok
  \sketch
  Starting from \cref{def:pet-ch3-second-fundamental-form}, metric compatibility
  and torsion-freeness give, for tangent $X,Y$,
  \[
   g(\nabla_XN,Y)=-g(N,\nabla_XY)
   =-g(N,\nabla_YX)=g(\nabla_YN,X),
  \]
  because $[X,Y]$ is tangent. Also
  $g(\nabla_XN,N)=\tfrac12Xg(N,N)=0$, so $\nabla_XN$ is tangent.
\end{proof}

\begin{proposition}[Level-set normals and Hessians]
  \label{prop:pet-ch3-normal-hessian-relation}
  \lean{PetersenLib.normal_hessian_relation}
  \source{petersen:page-0108,petersen:page-0109}
  \uses{def:pet-ch3-second-fundamental-form, def:pet-ch3-hessian-square-notation}
  \leanok
  \dcref{ch3:3.2.1}
  Let $H\subset f^{-1}(a)$ be open and consist entirely of regular points of $f$.
  Then:
  \begin{enumerate}
    \item[(1)] $N=\nabla f/|\nabla f|$ is a unit normal to $H$;
    \item[(2)] $\Pi(X,Y)=\dfrac1{|\nabla f|}\operatorname{Hess}f(X,Y)$ for
    $X,Y\in TH$;
    \item[(3)] $\operatorname{Hess}f(\nabla f,X)=\tfrac12D_X|\nabla f|^2$ for
    $X\in TM$.
  \end{enumerate}
\end{proposition}
\begin{proof}
  \uses{prop:pet-ch3-normal-hessian-relation}
  \leanok
  \sketch
  Tangency to a level set gives $g(\nabla f,X)=df(X)=0$, proving (1).
  Differentiating $N=\nabla f/|\nabla f|$ and pairing with tangent $Y$ removes the
  derivative of $|\nabla f|^{-1}$ and proves (2). Finally,
  \[
   \operatorname{Hess}f(\nabla f,X)
   =g(\nabla_X\nabla f,\nabla f)
   =\tfrac12X|\nabla f|^2,
  \]
  which proves (3).
\end{proof}

\begin{theorem}[Radial curvature equation]
  \label{thm:pet-ch3-radial-curvature-equation}
  \lean{PetersenLib.radialCurvatureEquation}
  \source{petersen:page-0110}
  \uses{def:pet-ch3-hessian-square-notation, def:pet-ch3-curvature-tensor}
  \leanok
  \dcref{ch3:3.2.2}
  With $S,\operatorname{Hess}^2f$ as in \cref{def:pet-ch3-hessian-square-notation}:
  \begin{align}
    (\nabla_{\nabla f}S)(X)+S^2(X)-\nabla_X(S(\nabla f))&=-R(X,\nabla f)\nabla f,\\
    \nabla_{\nabla f}\operatorname{Hess}f+\operatorname{Hess}^2f-\operatorname{Hess}\Big(\tfrac12|\nabla
    f|^2\Big)&=-R(\cdot,\nabla f,\nabla f,\cdot),\\
    \mathcal L_{\nabla f}\operatorname{Hess}f-\operatorname{Hess}^2f-\operatorname{Hess}\Big(\tfrac12|\nabla
    f|^2\Big)&=-R(\cdot,\nabla f,\nabla f,\cdot).
  \end{align}
\end{theorem}
\begin{proof}
  \uses{thm:pet-ch3-radial-curvature-equation}
  \leanok
  \sketch
  Set $S(X)=\nabla_X\nabla f$ and expand the curvature definition. The product
  rule for $\nabla S$ gives
  \[
   -R(X,\nabla f)\nabla f
   =(\nabla_{\nabla f}S)(X)+S^2(X)-\nabla_X(S(\nabla f)).
  \]
  Lowering an index and using
  $S(\nabla f)=\nabla(\tfrac12|\nabla f|^2)$ gives the second formula. Since
  \[
   \mathcal L_{\nabla f}\operatorname{Hess}f
   =\nabla_{\nabla f}\operatorname{Hess}f+2\operatorname{Hess}^2f,
  \]
  substitution gives the third.
\end{proof}

\begin{remark}[Tangential and normal decomposition]
  \label{rem:pet-ch3-tangential-normal-decomposition}
  \lean{PetersenLib.tangentialNormalDecomposition}
  \source{petersen:page-0110}
  \uses{thm:pet-ch3-radial-curvature-equation}
  \leanok
  \dcref{ch3:3.2.3}
  The third formula of \cref{thm:pet-ch3-radial-curvature-equation} computes a
  suitable curvature using only gradients of functions and Lie derivatives — no
  covariant derivatives are needed. The next two fundamental equations, known as
  the \emph{Gauss equation} and the \emph{Peterson--Codazzi--Mainardi equation},
  use the notation
  $X=X^\top+X^\perp=X-g(X,N)N+g(X,N)N$
  decomposing a vector into parts tangential and normal to a hypersurface $H$
  with normal $N$.
\end{remark}

\begin{theorem}[Tangential Gauss equation]
  \label{thm:pet-ch3-tangential-curvature-equation}
  \lean{PetersenLib.tangentialCurvatureEquation}
  \source{petersen:page-0111}
  \uses{def:pet-ch3-second-fundamental-form, rem:pet-ch3-tangential-normal-decomposition}
  \leanok
  \dcref{ch3:3.2.4}
  For $X,Y,Z,W$ tangent to a hypersurface $H$ with induced metric $g_H$ and
  induced curvature $R^H$,
  \[
    g(R(X,Y)Z,W)=g_H\big(R^H(X,Y)Z,W\big)-\Pi(X,W)\Pi(Y,Z)+\Pi(X,Z)\Pi(Y,W).
  \]
\end{theorem}
\begin{proof}
  \uses{thm:pet-ch3-tangential-curvature-equation}
  \sketch
  Using \cref{def:pet-ch3-second-fundamental-form}, the induced connection is
  $\nabla^H_XY=(\nabla_XY)^\top=\nabla_XY+\Pi(X,Y)N$. Substitute
  $\nabla_XY=\nabla^H_XY-\Pi(X,Y)N$ into $R(X,Y)Z$ and separate tangent and normal
  parts. Pairing the tangent part with $W$ gives
  \[
   g(R(X,Y)Z,W)=g_H(R^H(X,Y)Z,W)
   -\Pi(X,W)\Pi(Y,Z)+\Pi(X,Z)\Pi(Y,W).
  \]
\end{proof}

\begin{theorem}[Normal Codazzi--Mainardi equation]
  \label{thm:pet-ch3-normal-curvature-equation}
  \lean{PetersenLib.normalCurvatureEquation}
  \source{petersen:page-0111}
  \uses{thm:pet-ch3-tangential-curvature-equation}
  \leanok
  \dcref{ch3:3.2.5}
  For $X,Y,Z$ tangent to a hypersurface $H$,
  \[
    g(R(X,Y)Z,N)=-(\nabla_X\Pi)(Y,Z)+(\nabla_Y\Pi)(X,Z).
  \]
\end{theorem}
\begin{proof}
  \uses{thm:pet-ch3-normal-curvature-equation}
  \sketch
  In the tangent-normal expansion used for the Gauss equation, the normal
  coefficient is
  $-(\nabla_X\Pi)(Y,Z)+(\nabla_Y\Pi)(X,Z)$. Pairing with the unit normal $N$
  gives the formula.
\end{proof}

\begin{remark}[curvature equations by dimension]
  \label{rem:pet-ch3-curvature-equations-by-dimension}
  \lean{PetersenLib.curvatureEquationsByDimension}
  \source{petersen:page-0112}
  \uses{thm:pet-ch3-radial-curvature-equation, thm:pet-ch3-tangential-curvature-equation, thm:pet-ch3-normal-curvature-equation}
  \leanok
  Together the three fundamental equations compute curvature by induction on
  dimension: knowing curvature on $H$ and the shape operator $S$ determines
  curvature on $M$ at points of $H$. If $\dim M=1$ then $\dim H=0$, matching
  $R\equiv0$ on all $1$-dimensional spaces. If $\dim M=2$ then $\dim H=1$, so
  $R^H\equiv0$ and $X,Y,Z$ are proportional; only the radial curvature equation
  is relevant. If $\dim M=3$ then $\dim H=2$: writing an orthonormal frame
  $\{X,Y,N\}$ with $X\perp Y$, the full curvature tensor depends on
  $g(R(X,N)N,Y)$, $g(R(X,N)N,X)$, $g(R(Y,N)N,Y)$ (radial equation),
  $g(R(X,Y)Y,X)$ (tangential equation), and $g(R(X,Y)Y,N)$, $g(R(Y,X)X,N)$
  (normal equation).
\end{remark}

\begin{corollary}[Gauss's Theorema Egregium]
  \label{cor:pet-ch3-gauss-egregium}
  \lean{PetersenLib.gaussTheoremaEgregium}
  \source{petersen:page-0112}
  \uses{thm:pet-ch3-tangential-curvature-equation, ex:pet-ch3-euclidean-flat}
  \leanok
  When $M^3=\mathbb R^3$ (so $R\equiv0$) and $E_1,E_2$ is an orthonormal basis of
  $T_pH$ for a surface $H\subset\mathbb R^3$,
  \[
    \sec(T_pH)=R^H(E_1,E_2,E_2,E_1)=\Pi(E_1,E_1)\Pi(E_2,E_2)-\Pi(E_1,E_2)^2=\det[\Pi],
  \]
  i.e.\ the extrinsic quantity $\det[\Pi]$ equals the intrinsic sectional
  curvature $\sec(T_pH)$.
\end{corollary}
\begin{proof}
  \uses{cor:pet-ch3-gauss-egregium}
  \sketch
  Apply the Gauss equation to $X=W=E_1$, $Y=Z=E_2$. The ambient curvature
  vanishes by \cref{ex:pet-ch3-euclidean-flat}, and orthonormality identifies the intrinsic curvature term with
  $\sec(T_pH)$, leaving $\sec(T_pH)=\det[\Pi]$.
\end{proof}

\subsection{Distance Functions}

Distance functions specialize the level-set equations by imposing
$|\nabla r|=1$.

\begin{definition}[distance function]
  \label{def:pet-ch3-distance-function}
  \lean{PetersenLib.IsDistanceFunction}
  \source{petersen:page-0113}
  \leanok
  For $O\subset(M,g)$ open, $r:O\to\mathbb R$ is a \emph{distance function} if
  $|\nabla r|\equiv1$ on $O$, i.e.\ a solution of the \emph{Hamilton--Jacobi} (or
  \emph{eikonal}) equation $|\nabla r|^2=1$.
\end{definition}

\begin{example}[Euclidean point-distance functions]
  \label{ex:pet-ch3-distance-to-points}
  \lean{PetersenLib.distanceFunction_points}
  \source{petersen:page-0113}
  \uses{def:pet-ch3-distance-function}
  \leanok
  \dcref{ch3:3.2.6}
  On $(\mathbb R^n,g_{\mathrm{can}})$, $r(x)=|x-y|$ is smooth on $\mathbb
  R^n\setminus\{y\}$ with $|\nabla r|\equiv1$. For two points $y,z$,
  $r(x)=\operatorname{dist}(x,\{y,z\})=\min\{|x-y|,|x-z|\}$ is smooth away from
  $\{y,z\}$ and the
  hyperplane equidistant from $y,z$.
\end{example}

\begin{example}[Distance to a submanifold]
  \label{ex:pet-ch3-distance-to-submanifold}
  \lean{PetersenLib.distanceFunction_submanifold}
  \source{petersen:page-0113}
  \uses{def:pet-ch3-distance-function}
  \leanok
  \dcref{ch3:3.2.7}
  % Lean: formalized for a linear subspace; the curved case requires tubular neighborhoods.
  For a submanifold $H\subset\mathbb R^n$,
  $r(x)=\operatorname{dist}(x,H)=\inf\{|x-y|\mid y\in H\}$ is a distance function
  on some open $O$. When $H$ is an orientable hypersurface
  with unit normal $N$, coordinatize $\mathbb R^n$ by $x=tN+y$, $t\in\mathbb
  R,y\in H$; there is $\varepsilon:H\to(0,\infty)$ so that
  $O=\{tN+y\mid y\in H,|t|<\varepsilon(y)\}$ has unique coordinates $(t,y)$, and
  $r(x)=t$ (the \emph{signed distance}) or
  $f(x)=\operatorname{dist}(x,H)=|t|$ are distance functions on $O$, resp.\
  $O\setminus H$. More generally, on $I\times H$ with a metric
  $dr^2+g_r$ ($g_r$ depending on $r$), the projection $I\times H\to I$ is a
  distance function; special cases include rotationally symmetric metrics,
  doubly warped products, and submersion metrics on $I\times S^{2n-1}$.
\end{example}

\begin{lemma}[Distance functions as Riemannian submersions]
  \label{lem:pet-ch3-distance-function-submersion}
  \lean{PetersenLib.distanceFunction_iff_riemannianSubmersion}
  \source{petersen:page-0114}
  \uses{def:pet-ch3-distance-function}
  \leanok
  \dcref{ch3:3.2.8}
  For $O$ open in a Riemannian manifold, $r:O\to I\subset\mathbb R$ is a distance
  function if and only if it is a Riemannian submersion.
\end{lemma}
\begin{proof}
  \uses{lem:pet-ch3-distance-function-submersion}
  \leanok
  \sketch
  Since
  $Dr(v)=g(\nabla r,v)\partial_t$, the horizontal subspace
  $(\ker Dr)^\perp$ is spanned by $\nabla r$. On $v=\alpha\nabla r$,
  \[
   |v|=|\alpha||\nabla r|,\qquad |Dr(v)|=|\alpha||\nabla r|^2.
  \]
  Thus $Dr$ is an isometry on horizontal vectors exactly when $|\nabla r|=1$.
\end{proof}

\begin{definition}[distance-function notation: $\partial_r,O_r,g_r,S$]
  \label{def:pet-ch3-distance-function-notation}
  \lean{PetersenLib.distanceFunctionNotation}
  \source{petersen:page-0114}
  \uses{lem:pet-ch3-distance-function-submersion, def:pet-ch3-second-fundamental-form}
  \leanok
  Fix a distance function $r:O\to\mathbb R$. Write $\partial_r=\nabla r$; the
  level sets are $O_r=\{x\in O\mid r(x)=r\}$ with induced metric $g_r$, and
  $\nabla',R'$ denote the connection and curvature of $(O_r,g_r)$. Since
  $|\nabla r|=1$, $\operatorname{Hess}r=\Pi$, and $S$ denotes the common
  $(1,1)$-tensor for both: $S=\nabla\partial_r$ measures how the induced metric
  on $O_r$ changes, by measuring how the unit normal to $O_r$ changes.
\end{definition}

\begin{example}[Zero shape operator and hyperplanes]
  \label{ex:pet-ch3-zero-shape-operator-hyperplane}
  \lean{PetersenLib.zeroShapeOperator_iff_hyperplane}
  \source{petersen:page-0114,petersen:page-0115}
  \uses{def:pet-ch3-distance-function-notation}
  \dcref{ch3:3.2.9}
  Let $H\subset\mathbb R^n$ be an orientable hypersurface with unit normal $N$
  and shape operator $S(v)=\nabla_vN$. If $S=0$ on $H$, $N$ is constant on $H$
  and $H$ is open in the hyperplane $\{x+p\mid x\cdot N_p=0\}$ for fixed $p\in
  H$. Explicitly, for the isometric immersion
  $(x^1,\dots,x^{n-1})\mapsto(c(x^1),x^2,\dots,x^{n-1})$ of a unit-speed curve $c$
  into $\mathbb R^n$, the unit normal is $N=(-\dot c^2(x^1),\dot c^1(x^1),0,\dots,0)$,
  and $S\equiv0$ iff $\ddot c^1\equiv\ddot c^2\equiv0$ iff $c$ is a straight line
  iff $H$ is open in a hyperplane.
\end{example}
\begin{proof}
  \uses{ex:pet-ch3-zero-shape-operator-hyperplane}
  \sketch
  For the displayed immersion,
  \[
   \nabla N=(-\ddot c^2\partial_1+\ddot c^1\partial_2)\,dx^1.
  \]
  Hence $S=0$ exactly when $c$ is a straight line. In general, $S=0$ makes the
  unit normal constant on each connected component, so the hypersurface lies
  openly in the corresponding affine hyperplane; the converse is immediate.
\end{proof}

\begin{remark}[intrinsic versus extrinsic geometry]
  \label{rem:pet-ch3-intrinsic-extrinsic}
  \lean{PetersenLib.intrinsicExtrinsicGeometry}
  \source{petersen:page-0115}
  \uses{ex:pet-ch3-zero-shape-operator-hyperplane}
  \leanok
  The curvature tensor depends only on $(M,g)$, whereas the shape operator also
  depends on a chosen isometric immersion
  $(M,g)\to(\bar M,\bar g)$.
\end{remark}

\subsection{The Curvature Equations for Distance Functions}

The radial equation becomes a Riccati equation for the level-set shape
operator.

\begin{corollary}[Radial equation for distance functions]
  \label{cor:pet-ch3-radial-equation-distance-function}
  \lean{PetersenLib.radialEquation_distanceFunction}
  \source{petersen:page-0115}
  \uses{def:pet-ch3-distance-function-notation, thm:pet-ch3-radial-curvature-equation, prop:pet-ch3-normal-hessian-relation}
  \leanok
  \dcref{ch3:3.2.10}
  For a distance function $r:O\to\mathbb R$, $\nabla_{\partial_r}\partial_r=0$ and
  $\nabla_{\partial_r}S+S^2=-R_{\partial_r}$, where $R_{\partial_r}(X)=R(X,\partial_r)\partial_r$.
\end{corollary}
\begin{proof}
  \uses{cor:pet-ch3-radial-equation-distance-function}
  \leanok
  \sketch
  Because $|\nabla r|=1$, \cref{prop:pet-ch3-normal-hessian-relation} gives
  $\nabla_{\partial_r}\partial_r=0$. Substitute $f=r$ in the radial curvature
  equation; the term involving $S(\partial_r)$ vanishes, leaving
  $\nabla_{\partial_r}S+S^2=-R_{\partial_r}$.
\end{proof}

\begin{proposition}[Curvature equations for distance functions]
  \label{prop:pet-ch3-distance-function-curvature-equations}
  \lean{PetersenLib.distanceFunction_curvatureEquations}
  \source{petersen:page-0115,petersen:page-0116}
  \uses{cor:pet-ch3-radial-equation-distance-function, thm:pet-ch3-radial-curvature-equation}
  \leanok
  \dcref{ch3:3.2.11}
  For a distance function $r:(O,g)\to\mathbb R$ with $\partial_r=\nabla r$:
  \begin{enumerate}
    \item[(1)] $\mathcal L_{\partial_r}g=2\operatorname{Hess}r$;
    \item[(2)] $(\nabla_{\partial_r}\operatorname{Hess}r)(X,Y)+\operatorname{Hess}^2r(X,Y)=-R(X,\partial_r,\partial_r,Y)$;
    \item[(3)] $(\mathcal L_{\partial_r}\operatorname{Hess}r)(X,Y)-\operatorname{Hess}^2r(X,Y)=-R(X,\partial_r,\partial_r,Y)$.
  \end{enumerate}
\end{proposition}
\begin{proof}
  \uses{prop:pet-ch3-distance-function-curvature-equations}
  \leanok
  \sketch
  The identity
  $\mathcal L_{\nabla f}g=2\operatorname{Hess}f$ gives (1). Apply the two
  $(0,2)$ radial curvature equations to $f=r$; since $|\nabla r|=1$,
  $\operatorname{Hess}(\tfrac12|\nabla r|^2)=0$, yielding (2) and (3).
\end{proof}

\subsection{Jacobi Fields}

Fields invariant under the radial flow convert the tensor evolution equations
into ordinary differential equations along radial curves.

\begin{definition}[Jacobi field]
  \label{def:pet-ch3-jacobi-field}
  \lean{PetersenLib.IsJacobiField}
  \source{petersen:page-0116}
  \uses{def:pet-ch3-distance-function-notation}
  \leanok
  For a smooth distance function $r$, a \emph{Jacobi field} is a smooth vector
  field $J$ independent of $r$, i.e.\ satisfying the \emph{Jacobi equation}
  $\mathcal L_{\partial_r}J=0$.
\end{definition}

\begin{remark}[construction of Jacobi fields]
  \label{rem:pet-ch3-jacobi-field-construction}
  \lean{PetersenLib.jacobiField_construction}
  \source{petersen:page-0116}
  \uses{def:pet-ch3-jacobi-field}
  \leanok
  \sketch
  In coordinates $(r,x^2,\dots,x^n)$ adapted to $r$, write
  $J=J^r\partial_r+J^i\partial_i$. Then
  $\mathcal L_{\partial_r}J=\partial_r(J^r)\partial_r+\partial_r(J^i)\partial_i$, so the Jacobi
  equation is equivalent to $\partial_rJ^r=\partial_rJ^i=0$. Hence values on a
  hypersurface transverse to $\partial_r$ determine $J$, and the coordinate
  vector fields are Jacobi fields.
\end{remark}

\begin{lemma}[first-order form of the Jacobi equation]
  \label{lem:pet-ch3-jacobi-field-equivalent-equation}
  \lean{PetersenLib.jacobiField_equivalentEquation}
  \source{petersen:page-0116}
  \uses{def:pet-ch3-jacobi-field}
  \leanok
  $\mathcal L_{\partial_r}J=0$ is equivalent to $\nabla_{\partial_r}J=S(J)$, and
  consequently $\operatorname{Hess}r(J,J)=g(\nabla_{\partial_r}J,J)=\tfrac12\partial_rg(J,J)$.
\end{lemma}
\begin{proof}
  \uses{lem:pet-ch3-jacobi-field-equivalent-equation}
  \leanok
  \sketch
  Torsion-freeness gives
  $\mathcal L_{\partial_r}J=[\partial_r,J]
  =\nabla_{\partial_r}J-\nabla_J\partial_r
  =\nabla_{\partial_r}J-S(J)$.
  Pairing the resulting equation with $J$ proves the Hessian and length identity.
\end{proof}

\begin{proposition}[second-order Jacobi equation]
  \label{prop:pet-ch3-jacobi-second-order-equation}
  \lean{PetersenLib.jacobiField_secondOrderEquation}
  \source{petersen:page-0117}
  \uses{lem:pet-ch3-jacobi-field-equivalent-equation, cor:pet-ch3-radial-equation-distance-function, def:pet-ch3-directional-curvature-operator}
  \leanok
  A Jacobi field $J$ satisfies $\nabla_{\partial_r}\nabla_{\partial_r}J=-R(J,\partial_r)\partial_r$.
\end{proposition}
\begin{proof}
  \uses{prop:pet-ch3-jacobi-second-order-equation}
  \leanok
  \sketch
  Differentiate $\nabla_{\partial_r}J=S(J)$ along $\partial_r$:
  \[
   \nabla_{\partial_r}^2J=(\nabla_{\partial_r}S)(J)+S^2(J).
  \]
  The radial equation \cref{cor:pet-ch3-radial-equation-distance-function}
  identifies the right-hand side with
  $-R(J,\partial_r)\partial_r=-R_{\partial_r}(J)$ in the notation of
  \cref{def:pet-ch3-directional-curvature-operator}.
\end{proof}

\begin{proposition}[fundamental equations on Jacobi fields]
  \label{prop:pet-ch3-jacobi-field-curvature-equations}
  \lean{PetersenLib.jacobiField_curvatureEquations}
  \source{petersen:page-0117}
  \uses{def:pet-ch3-jacobi-field, prop:pet-ch3-distance-function-curvature-equations}
  \leanok
  For Jacobi fields $J_1,J_2$, equations (1) and (3) of
  \cref{prop:pet-ch3-distance-function-curvature-equations} become
  \[
    \partial_rg(J_1,J_2)=2\operatorname{Hess}r(J_1,J_2),\qquad
    \partial_r\big(\operatorname{Hess}r(J_1,J_2)\big)-\operatorname{Hess}^2r(J_1,J_2)=-R(J_1,\partial_r,\partial_r,J_2).
  \]
\end{proposition}
\begin{proof}
  \uses{prop:pet-ch3-jacobi-field-curvature-equations}
  \leanok
  \sketch
  Evaluate the Lie-derivative equations for $g$ and $\operatorname{Hess}r$ on
  $J_1,J_2$. The terms differentiating the fields vanish
  because $\mathcal L_{\partial_r}J_i=0$, leaving the two ordinary
  $\partial_r$-derivatives in the statement.
\end{proof}

\begin{remark}[reduction on a product neighborhood]
  \label{rem:pet-ch3-jacobi-field-product-neighborhood}
  \lean{PetersenLib.jacobiField_productNeighborhood}
  \source{petersen:page-0117}
  \uses{prop:pet-ch3-jacobi-field-curvature-equations}
  \leanok
  \sketch
  On a product neighborhood $\Omega=(a,b)\times H\subset M$ with $r(t,z)=t$, the
  metric is $g=dr^2+g_r$, a one-parameter family of metrics on $H$; a vector
  field on $H$ extends uniquely to a Jacobi field on $\Omega$. Since
  $\operatorname{Hess}r(\partial_r,J)=g(\nabla_{\partial_r}\partial_r,J)=0$ and
  $g_r(\partial_r,J)=0$, restricting to $H$ gives $\partial_rg=\partial_rg_r=2\operatorname{Hess}r$,
  and the fundamental equations become, on $H$,
  \[
    \partial_rg_r=2\operatorname{Hess}r,\qquad
    \partial_r\operatorname{Hess}r-\operatorname{Hess}^2r=-R(\cdot,\partial_r,\partial_r,\cdot).
  \]
\end{remark}

\begin{lemma}[the radial curvature term via sectional curvature]
  \label{lem:pet-ch3-jacobi-field-sectional-curvature-formula}
  \lean{PetersenLib.jacobiField_sectionalCurvatureFormula}
  \source{petersen:page-0118}
  \uses{rem:pet-ch3-jacobi-field-product-neighborhood, def:pet-ch3-sectional-curvature}
  \leanok
  For $X$ tangent to $H$ (so $g(X,\partial_r)=0$),
  $R(X,\partial_r,\partial_r,X)=\sec(X,\partial_r)g_r(X,X)$; hence for a Jacobi
  field $J$ on $H$,
  \[
    \partial_r\big(\operatorname{Hess}r(J,J)\big)-\operatorname{Hess}^2r(J,J)=-\sec(J,\partial_r)g_r(J,J).
  \]
  This still couples the Hessian equation to the metric $g_r$.
\end{lemma}
\begin{proof}
  \uses{lem:pet-ch3-jacobi-field-sectional-curvature-formula}
  \leanok
  \sketch
  Orthogonality to the unit field $\partial_r$ gives
  \[
   R(X,\partial_r,\partial_r,X)
   =\sec(X,\partial_r)g(X,X).
  \]
  For tangent $X$, $g(X,X)=g_r(X,X)$; substituting $X=J$ in the
  product-neighborhood Hessian equation of
  \cref{rem:pet-ch3-jacobi-field-product-neighborhood} proves the second formula.
\end{proof}

\begin{example}[Gaussian curvature in distance coordinates]
  \label{ex:pet-ch3-two-dim-sectional-curvature-formula}
  \lean{PetersenLib.twoDimensional_sectionalCurvature_formula}
  \source{petersen:page-0118,petersen:page-0119}
  \uses{lem:pet-ch3-jacobi-field-sectional-curvature-formula}
  \leanok
  \dcref{ch3:3.2.12}
  If $\dim M=2$, write $g=dr^2+\rho^2(r,\theta)d\theta^2$ with $\theta$
  coordinatizing the level sets of $r$. Then $\partial_\theta$ is a Jacobi field
  of length $\rho$, and $\sec(T_pM)=-\partial_r^2\rho/\rho$.
\end{example}
\begin{proof}
  \uses{ex:pet-ch3-two-dim-sectional-curvature-formula}
  \leanok
  \sketch
  By \cref{rem:pet-ch3-jacobi-field-construction},
  $J=\partial_\theta$ is a Jacobi field. The equations in
  \cref{prop:pet-ch3-jacobi-field-curvature-equations} give
  $\operatorname{Hess}r(J,J)=\rho\partial_r\rho$ and
  $S(J)=(\partial_r\rho/\rho)J$. Hence
  \[
   -\sec(J,\partial_r)\rho^2
   =\partial_r(\rho\partial_r\rho)-(\partial_r\rho)^2
   =\rho\partial_r^2\rho,
  \]
  where \cref{lem:pet-ch3-jacobi-field-sectional-curvature-formula} supplies the
  left-hand side. This yields the claimed formula.
\end{proof}

\subsection{Parallel Fields}

Parallel transport gives a second class of test fields for the radial equation.

\begin{definition}[parallel field]
  \label{def:pet-ch3-parallel-field}
  \lean{PetersenLib.IsParallelField}
  \source{petersen:page-0119}
  \uses{def:pet-ch3-distance-function-notation}
  \leanok
  For a smooth distance function $r$, a \emph{parallel field} is a vector field
  $X$ with $\nabla_{\partial_r}X=0$. Parallel fields are almost never Jacobi
  fields.
\end{definition}

\begin{proposition}[curvature equation on parallel fields]
  \label{prop:pet-ch3-parallel-field-curvature-equation}
  \lean{PetersenLib.parallelField_curvatureEquation}
  \source{petersen:page-0119}
  \uses{def:pet-ch3-parallel-field, prop:pet-ch3-distance-function-curvature-equations, lem:pet-ch3-jacobi-field-sectional-curvature-formula}
  \leanok
  For parallel fields $X,Y$, $\partial_rg(X,Y)=0$; equation (2) of
  \cref{prop:pet-ch3-distance-function-curvature-equations} becomes
  $\partial_r(\operatorname{Hess}r(X,Y))+\operatorname{Hess}^2r(X,Y)=-R(X,\partial_r,\partial_r,Y)$,
  and on a unit parallel field $X$ orthogonal to $\partial_r$,
  \[
    \partial_r\big(\operatorname{Hess}r(X,X)\big)+\operatorname{Hess}^2r(X,X)=-\sec(X,\partial_r),
  \]
  decoupled from the metric since $g_r(X,X)$ is then constant in $r$.
\end{proposition}
\begin{proof}
  \uses{prop:pet-ch3-parallel-field-curvature-equation}
  \leanok
  \sketch
  Metric compatibility gives $\partial_rg(X,Y)=0$ for parallel fields. Evaluate
  the covariant radial Hessian equation on $X,Y$; no field-derivative terms remain.
  For unit $X\perp\partial_r$, the sectional-curvature identity gives
  $R(X,\partial_r,\partial_r,X)=\sec(X,\partial_r)$.
\end{proof}

\subsection{Conjugate Points}

Degeneration of the radial Hessian or intersection of radial trajectories leads
to the following two notions.

\begin{definition}[conjugate point, focal point]
  \label{def:pet-ch3-conjugate-focal-point}
  \lean{PetersenLib.ConjugatePoint}
  \source{petersen:page-0120}
  \uses{prop:pet-ch3-distance-function-curvature-equations, prop:pet-ch3-parallel-field-curvature-equation}
  \leanok
  Along an integral curve of $\partial_r$ with bounded curvature, if
  $\operatorname{Hess}r$ blows up, equation (2) forces it to be decreasing as $r$
  increases, so it can only tend to $-\infty$; by equation (1), the only possible
  degeneration along the curve is that the metric $g$ stops being positive
  definite. When this happens $r$ develops a \emph{conjugate} or \emph{focal
  point} along the curve. A \emph{conjugate point} occurs when
  $\operatorname{Hess}r$ becomes undefined while solving its differential
  equation; a \emph{focal point} occurs when integral curves of $\nabla r$ meet.
  The two frequently coincide but need not: there are metrics with conjugate
  points that are not focal points.
\end{definition}

\begin{remark}[conjugate and focal loci of an ellipse]
  \label{rem:pet-ch3-conjugate-vs-focal-example}
  \lean{PetersenLib.conjugateVsFocal_ellipse}
  \source{petersen:page-0120}
  \uses{def:pet-ch3-conjugate-focal-point}
  \leanok
  % Lean: formalized as plane geometry; no abstract focal-point API is used.
  \sketch
  For $c(t)=(a\cos t,b\sin t)$ with $0<b<a$, the conjugate locus along the
  normals is the evolute

  $e(t)=\bigl(\tfrac{a^2-b^2}{a}\cos^3t,
  \tfrac{b^2-a^2}{b}\sin^3t\bigr)$.
  The normals at $t$ and $-t$ meet at
  $\bigl(\tfrac{a^2-b^2}{a}\cos t,0\bigr)$, so the focal locus is the major-axis
  segment with endpoints $\pm\tfrac{a^2-b^2}{a}$. These are the cusps of the
  evolute, not the conic foci $\pm\sqrt{a^2-b^2}$, since
  \[
    \frac{a^2-b^2}{a}<\sqrt{a^2-b^2}\qquad\text{for every }0<b<a
  \]
  Thus the focal locus is strictly contained in the segment joining the conic
  foci.
\end{remark}
\begin{remark}[sign consequences of radial Hessian evolution]
  \label{rem:pet-ch3-external-force-internal-reaction}
  \lean{PetersenLib.externalForceInternalReaction}
  \source{petersen:page-0120}
  \uses{def:pet-ch3-conjugate-focal-point, prop:pet-ch3-distance-function-curvature-equations}
  \leanok
  \sketch
  Equations (2) and (3) of
  \cref{prop:pet-ch3-distance-function-curvature-equations} can be written as
  \[
    (\nabla_{\partial_r}\operatorname{Hess}r)(X,X)=-R(X,\partial_r,\partial_r,X)-\operatorname{Hess}^2r(X,X),\qquad
    (\mathcal L_{\partial_r}\operatorname{Hess}r)(X,X)=-R(X,\partial_r,\partial_r,X)+\operatorname{Hess}^2r(X,X),
  \]
  Hence, if $\sec\le0$, $\mathcal L_{\partial_r}\operatorname{Hess}r\ge0$: positive
  $\operatorname{Hess}r$ stays positive. If $\sec\ge0$,
  $\nabla_{\partial_r}\operatorname{Hess}r\le0$: nonpositive
  $\operatorname{Hess}r$ stays nonpositive.
\end{remark}

\section{Exercises}

\begin{remark}[Jet dependence of induced curvature]
  \label{rem:pet-ch3-ex-1}
  \lean{PetersenLib.exercise3_4_1}
  \source{petersen:page-0121}
  \uses{def:pet-ch3-curvature-tensor}
  \leanok
  \dcref{ch3:3.4.1}
  % Lean: formalized as invariance of induced curvature under equality of first and second jets.
  Let $M$ be an $n$-dimensional submanifold of $\mathbb R^{n+m}$ with the induced
  metric, given locally by a parametrization $u^s(x^1,\dots,x^n)$,
  $s=1,\dots,n+m$. Show that in these coordinates $R^l_{ijk}$ depends only on the
  first and second partials of $u^s$.
\end{remark}

\begin{remark}[Equivalent gradient-length conditions]
  \label{rem:pet-ch3-ex-2}
  \lean{PetersenLib.exercise3_4_2}
  \source{petersen:page-0121}
  \uses{def:pet-ch3-hessian-square-notation}
  \leanok
  \dcref{ch3:3.4.2}
  For $f:(M,g)\to\mathbb R$ smooth on a connected Riemannian manifold, consider:
  (1) $|\nabla f|$ is constant; (2) $\nabla_{\nabla f}\nabla f=0$; (3) $|\nabla
  f|$ is constant on the level sets of $f$. Show $(1)\Leftrightarrow(2)\Rightarrow(3)$,
  and give an example showing the last implication is not an equivalence.
\end{remark}

\begin{remark}[Lie and covariant derivatives of the Hessian operator]
  \label{rem:pet-ch3-ex-3}
  \lean{PetersenLib.exercise3_4_3}
  \source{petersen:page-0121}
  \uses{def:pet-ch3-hessian-square-notation, thm:pet-ch3-radial-curvature-equation}
  \leanok
  \dcref{ch3:3.4.3}
  For $f$ with $S(X)=\nabla_X\nabla f$, show that
  $\mathcal L_{\nabla f}S=\nabla_{\nabla f}S$ and
  $(\mathcal L_{\nabla f}S)(X)+S^2(X)-\nabla_X(S(\nabla f))=-R_{\nabla f}(X)$.
  Reconcile this
  with \cref{thm:pet-ch3-radial-curvature-equation} for the $(0,2)$-version of
  the Hessian.
\end{remark}

\begin{remark}[Gauss and Codazzi equations for distance functions]
  \label{rem:pet-ch3-ex-4}
  \lean{PetersenLib.exercise3_4_4}
  \source{petersen:page-0121,petersen:page-0122}
  \uses{thm:pet-ch3-tangential-curvature-equation, thm:pet-ch3-normal-curvature-equation}
  \leanok
  \dcref{ch3:3.4.4}
  For a distance function $r=f:M\to\mathbb R$, show the tangential and mixed
  curvature equations can be written as
  \[
    (R(X,Y)Z)^\top=R_H(X,Y)Z-(S(X)\wedge S(Y))(Z),\qquad
    g(R(X,Y)Z,N)=-g((\nabla_XS)(Y),Z)+g((\nabla_YS)(X),Z),
  \]
  and $R(X,Y)N=(d^\nabla S)(X,Y)$.
\end{remark}

\begin{remark}[Bianchi identities in normal coordinates]
  \label{rem:pet-ch3-ex-5}
  \lean{PetersenLib.exercise3_4_5}
  \source{petersen:page-0122}
  \uses{prop:pet-ch3-curvature-symmetries}
  \leanok
  \dcref{ch3:3.4.5}
  Prove both Bianchi identities at a point $p\in M$ using a coordinate system
  where $\nabla_{\partial_i}\partial_j=0$ at $p$.
\end{remark}

\begin{remark}[Parallel curvature of space forms]
  \label{rem:pet-ch3-ex-6}
  \lean{PetersenLib.exercise3_4_6}
  \source{petersen:page-0122}
  \uses{def:pet-ch3-constant-curvature}
  \leanok
  \dcref{ch3:3.4.6}
  Show that a Riemannian manifold with constant curvature has parallel curvature
  tensor.
\end{remark}

\begin{remark}[Parallel Ricci and scalar curvature]
  \label{rem:pet-ch3-ex-7}
  \lean{PetersenLib.exercise3_4_7}
  \source{petersen:page-0122}
  \uses{def:pet-ch3-ricci-curvature, def:pet-ch3-scalar-curvature}
  \leanok
  \dcref{ch3:3.4.7}
  Show that a Riemannian manifold with parallel Ricci tensor has constant scalar
  curvature. (The converse fails, and a metric with parallel curvature tensor
  need not be Einstein.)
\end{remark}

\begin{remark}[Divergence of the curvature tensor]
  \label{rem:pet-ch3-ex-8}
  \lean{PetersenLib.exercise3_4_8}
  \source{petersen:page-0122}
  \uses{prop:pet-ch3-contracted-bianchi, notation:pet-ch3-divergence-adjoint}
  \leanok
  \dcref{ch3:3.4.8}
  With $R$ the $(0,4)$-curvature tensor and $\operatorname{Ric}$ the
  $(0,2)$-Ricci tensor, show
  $(\nabla^*R)(Z,X,Y)=(\nabla_X\operatorname{Ric})(Y,Z)-(\nabla_Y\operatorname{Ric})(X,Z)$.
  Conclude $\nabla^*R=0$ if $\nabla\operatorname{Ric}=0$, and that $\nabla^*R=0$
  iff the $(1,1)$-Ricci tensor satisfies $(\nabla_X\operatorname{Ric})(Y)=(\nabla_Y\operatorname{Ric})(X)$
  for all $X,Y$.
\end{remark}

\begin{remark}[Mixed curvature of product metrics]
  \label{rem:pet-ch3-ex-9}
  \lean{PetersenLib.exercise3_4_9}
  \source{petersen:page-0122}
  \uses{def:pet-ch3-sectional-curvature, rem:pet-ch1-ex-1}
  \leanok
  \dcref{ch3:3.4.9}
  % Lean: both the mixed covariant derivative and mixed sectional curvature statements are formalized.
  For Riemannian manifolds $(M,g_M),(N,g_N)$ with product metric
  $(M\times N,g_M+g_N)$, and $X$ a vector field on $M$, $Y$ on $N$ regarded on
  $M\times N$, show $\nabla_XY=0$; conclude that the mixed sectional curvature
  $\sec(X,Y)$ vanishes.
\end{remark}

\begin{remark}[Orthogonality characterization of constant curvature]
  \label{rem:pet-ch3-ex-10}
  \lean{PetersenLib.exercise3_4_10}
  \source{petersen:page-0122}
  \uses{def:pet-ch3-constant-curvature}
  \leanok
  \dcref{ch3:3.4.10}
  Show that $(M,g)$ has constant curvature at $p$ iff $R(v,w)z=0$ for all pairwise
  orthogonal $v,w,z\in T_pM$. Hint: a symmetric bilinear form $B$ on an inner
  product space with $B(v,w)=0$ whenever $v\perp w$ is a multiple of the inner
  product.
\end{remark}

\begin{remark}[General-codimension Gauss and Codazzi equations]
  \label{rem:pet-ch3-ex-11}
  \lean{PetersenLib.exercise3_4_11}
  \source{petersen:page-0122}
  \uses{thm:pet-ch3-tangential-curvature-equation}
  \leanok
  \dcref{ch3:3.4.11}
  % Lean: formalized via an abstract tangent-normal projection and the induced operator splitting.
  For $X,Y,Z$ tangent to $M$ (a submanifold with second fundamental form $T$ and
  normal connection $\nabla^\perp$), show
  \[
    \bar R(X,Y)Z=R^M(X,Y)Z+T_XT_YZ-T_YT_XZ+(\nabla^\perp_XT)_YZ-(\nabla^\perp_YT)_XZ,
  \]
  where $\bar R$ is the curvature of the ambient manifold and
  $(\nabla^\perp_XT)_YZ=\nabla^\perp_X(T_YZ)-T_{\nabla^M_XY}Z-T_Y\nabla^M_XZ$;
  the tangential parts recover the Gauss equations, the normal parts the
  Peterson--Codazzi--Mainardi equations. (The source prints $R^M$ on both sides;
  the left-hand curvature is the ambient one $\bar R$, not the induced $R^M$.)
\end{remark}

\begin{remark}[Ricci curvature of Euclidean hypersurfaces]
  \label{rem:pet-ch3-ex-12}
  \lean{PetersenLib.exercise3_4_12}
  \source{petersen:page-0123}
  \uses{def:pet-ch3-second-fundamental-form, def:pet-ch3-ricci-curvature}
  \leanok
  \dcref{ch3:3.4.12}
  For $H^{n-1}\subset\mathbb R^n$, show
  $\operatorname{Ric}^H=\operatorname{tr}(\Pi)\cdot\Pi-\Pi^2$.
\end{remark}

\begin{remark}[Totally geodesic hypersurfaces and constant curvature]
  \label{rem:pet-ch3-ex-13}
  \lean{PetersenLib.exercise3_4_13}
  \source{petersen:page-0123}
  \uses{def:pet-ch3-second-fundamental-form, def:pet-ch3-constant-curvature, rem:pet-ch3-ex-10}
  \leanok
  \dcref{ch3:3.4.13}
  % Lean: part (2) is formalized; part (1) requires space-form models.
  A hypersurface is \emph{totally geodesic} if its second fundamental form
  vanishes. (1) Show that the space forms $S^n_k$ have the property that any
  tangent vector is normal to a totally geodesic hypersurface. (2) Show that a
  Riemannian $n$-manifold, $n>2$, with this property has constant curvature.
  Hint: show $R(X,Y)Z=0$ for pairwise orthogonal $X,Y,Z$ and use
  \cref{rem:pet-ch3-ex-10}.
\end{remark}

\begin{remark}[Normal curvature and Ricci equations]
  \label{rem:pet-ch3-ex-14}
  \lean{PetersenLib.exercise3_4_14}
  \source{petersen:page-0123}
  \uses{def:pet-ch3-second-fundamental-form}
  \leanok
  \dcref{ch3:3.4.14}
  % Lean: the two skew-symmetries and Ricci equation are formalized for an abstract smooth orthogonal splitting.
  For tangent fields $X,Y$ and normal fields $V,W$ of a submanifold, define the
  normal curvature $R^\perp(X,Y,V,W)$. (1) Show $R^\perp$ is tensorial and
  skew-symmetric in $X,Y$ and in $V,W$. (2) Show
  $R^{\bar M}(X,Y,V,W)=R^\perp(X,Y,V,W)+g_M(T_XV,T_YW)-g_M(T_YV,T_XW)$ (the
  \emph{Ricci equations}).
\end{remark}

\begin{remark}[Curvature operator and Ricci curvature in dimension three]
  \label{rem:pet-ch3-ex-15}
  \lean{PetersenLib.exercise3_4_15}
  \source{petersen:page-0123}
  \uses{def:pet-ch3-curvature-operator, def:pet-ch3-ricci-curvature}
  \leanok
  \dcref{ch3:3.4.15}
  For a $3$-manifold, if the curvature operator in diagonal form is
  $\operatorname{diag}(\alpha,\beta,\gamma)$, show the Ricci curvature is diagonal
  with entries $\alpha+\beta,\beta+\gamma,\alpha+\gamma$, and that
  $\alpha,\beta,\gamma$ must be sectional curvatures.
\end{remark}

\begin{remark}[Einstein tensor]
  \label{rem:pet-ch3-ex-16}
  \lean{PetersenLib.exercise3_4_16}
  \source{petersen:page-0123,petersen:page-0124}
  \uses{def:pet-ch3-ricci-curvature, def:pet-ch3-scalar-curvature, notation:pet-ch3-divergence-adjoint}
  \leanok
  \dcref{ch3:3.4.16}
  For the $(0,2)$-tensor $T=\operatorname{Ric}+b\cdot\operatorname{scal}\cdot
  g+cg$, $b,c\in\mathbb R$: (1) show $\nabla^*T=0$ if $b=-\tfrac12$ — the tensor
  $G=\operatorname{Ric}-\tfrac{\operatorname{scal}}2g+cg$ is the \emph{Einstein
  tensor}, $c$ the \emph{cosmological constant}; (2) if $c=0$, show $G=0$ in
  dimension $2$; (3) for $n>2$, show $G=0$ implies the metric is Einstein; (4)
  for $n>2$, show $G=0$ and $c=0$ implies the metric is Ricci flat.
\end{remark}

\begin{remark}[Complex and isotropic sectional curvature]
  \label{rem:pet-ch3-ex-17}
  \lean{PetersenLib.exercise3_4_17}
  \source{petersen:page-0124,petersen:page-0125}
  \uses{def:pet-ch3-sectional-curvature}
  \leanok
  \dcref{ch3:3.4.17}
  % Lean: parts (1)--(3) are formalized; parts (4)--(5) are deferred.
  Let $T_{\mathbb C}M=TM\otimes\mathbb C$; for $v=v_1+iv_2\in T_{\mathbb C}M$
  ($v_1,v_2\in TM$), $\bar v=v_1-iv_2$, and any tensorial object on $TM$
  complexifies, e.g.\ for a $(1,1)$-tensor $S$,
  $S_{\mathbb C}(v_1+iv_2)=S(v_1)+iS(v_2)$. The complex-bilinear extension
  $g_{\mathbb C}$ of a Riemannian metric $g$ induces the Hermitian form
  $h(v,w)=g_{\mathbb C}(v,\bar w)
  =g(v_1,w_1)+g(v_2,w_2)+i(g(v_2,w_1)-g(v_1,w_2))$. A vector $v$
  is \emph{isotropic} if $g_{\mathbb C}(v,v)=0$, equivalently $h(v,\bar v)=0$, i.e.\
  $g(v_1,v_1)=g(v_2,v_2)$ and $g(v_1,v_2)=0$; more generally isotropic subspaces
  are those on which $g_{\mathbb C}$ vanishes. The \emph{complex sectional
  curvature} of $h$-orthonormal $v,w$ is $R_{\mathbb C}(v,w,\bar w,\bar v)$,
  called \emph{isotropic} when $v,w$ span an isotropic plane. (1) Show $v_1+iv_2$
  is isotropic if $v_1,v_2$ are orthogonal of equal length. (2) Show that if an
  isotropic plane is spanned by $h$-orthonormal isotropic $v=v_1+iv_2$,
  $w=w_1+iw_2$, then $v_1,v_2,w_1,w_2$ are pairwise orthogonal and each has
  squared norm $\tfrac12$ (the printed source calls them orthonormal). (3) Show
  $R_{\mathbb C}(v,w,\bar w,\bar v)$ is always real. (4) Show that if the metric
  is strictly quarter-pinched ($\sec\in(\tfrac14k,k)$, $k>0$), the complex
  sectional curvatures are positive. (5) Show the complex sectional curvatures
  are nonnegative (resp.\ positive) if the curvature operator is (resp.\
  strictly). Hint: compare
  $g(\mathfrak R(x\wedge u-y\wedge v),x\wedge u-y\wedge v)+g(\mathfrak R(x\wedge
  v+y\wedge u),x\wedge v+y\wedge u)$ to a suitable complex curvature.
\end{remark}

\begin{remark}[Curvature under metric scaling]
  \label{rem:pet-ch3-ex-18}
  \lean{PetersenLib.exercise3_4_18}
  \source{petersen:page-0125}
  \uses{def:pet-ch3-sectional-curvature, def:pet-ch3-scalar-curvature, def:pet-ch3-curvature-operator, def:pet-ch3-ricci-curvature}
  \leanok
  \dcref{ch3:3.4.18}
  % Lean: parts (1), (3), (4), and the sectional/scalar parts of (2) are formalized.
  Scaling $(M,g)$ to $(M,\lambda^2g)$: (1) show the connection and $(1,3)$-curvature
  tensor are unchanged; (2) $\sec,\operatorname{scal},\mathfrak R$ are all
  multiplied by $\lambda^{-2}$; (3) $\operatorname{Ric}$ as a $(1,1)$-tensor is
  multiplied by $\lambda^{-2}$; (4) $\operatorname{Ric}$ as a $(0,2)$-tensor is
  unchanged.
\end{remark}

\begin{remark}[Affine vector fields]
  \label{rem:pet-ch3-ex-19}
  \lean{PetersenLib.exercise3_4_19}
  \source{petersen:page-0125}
  \uses{def:pet-ch3-curvature-tensor}
  \leanok
  \dcref{ch3:3.4.19}
  Call $X$ an \emph{affine vector field} if $\mathcal L_X\nabla=0$. Show such a
  field satisfies $\nabla^2_{U,V}X=-R(X,U)V$.
\end{remark}

\begin{remark}[Integrability conditions for first-order systems]
  \label{rem:pet-ch3-ex-20}
  \lean{PetersenLib.exercise3_4_20}
  \source{petersen:page-0125,petersen:page-0126}
  \uses{prop:pet-ch3-curvature-coordinate-formula}
  \leanok
  \dcref{ch3:3.4.20}
  % Lean: necessity in part (1) is formalized; parts (2)--(4) require Frobenius-type existence.
  For $P^i_k(x,u)$, $x=(x^1,\dots,x^n)$, $u=(u^1,\dots,u^m)$, consider
  $\partial u^i/\partial x^k=P^i_k(x,u(x))$, $u(x_0)=u_0$. (1) Show
  $\partial^2u^i/\partial x^k\partial x^j=\partial P^i_k/\partial
  x^j+(\partial P^i_k/\partial u^l)P^l_j$, and conclude solvability requires the
  \emph{integrability conditions}
  $\partial P^i_k/\partial x^j+(\partial P^i_k/\partial u^l)P^l_j=\partial
  P^i_j/\partial x^k+(\partial P^i_j/\partial u^l)P^l_k$. (2) Conversely show
  the integrability conditions suffice (equivalent to the Frobenius theorem). (3)
  For the system $\partial U^i_j/\partial x^k=\sum_s\Gamma^s_{kj}U^i_s$ on a
  Riemannian $n$-manifold, show the integrability conditions are equivalent to
  $R^s_{klj}=0$. (4) Show a flat Riemannian manifold admits Cartesian
  coordinates, via $\partial u^i/\partial x^k=U^i_k$ with appropriate initial
  values, checking $g^{kl}(\partial u^i/\partial x^k)(\partial u^j/\partial x^l)$
  is constant using \cref{prop:pet-ch3-curvature-coordinate-formula}.
\end{remark}

\begin{remark}[Fundamental theorem of hypersurface theory]
  \label{rem:pet-ch3-ex-21}
  \lean{PetersenLib.exercise3_4_21}
  \source{petersen:page-0126,petersen:page-0127}
  \uses{prop:pet-ch3-curvature-coordinate-formula, thm:pet-ch3-tangential-curvature-equation, thm:pet-ch3-normal-curvature-equation}
  \leanok
  \dcref{ch3:3.4.21}
  % Lean: part (1) and the necessity direction of part (2) are formalized; the converse and existence claims require Frobenius. The converse implication in part (4) is formalized.
  For a Riemannian immersion $F=(u^1,\dots,u^{n+1}):M^n\hookrightarrow\mathbb
  R^{n+1}$, $U^i_k=\partial u^i/\partial x^k$: (1) show
  $\partial U^i_j/\partial x^k=\sum_s\Gamma^s_{kj}U^i_s-\Pi_{jk}N^i$ for a unit normal
  $N=N^i\partial/\partial u^i$ and $\Pi_{jk}=\Pi(\partial_j,\partial_k)=g(\nabla_jN,\partial_k)$;
  (2) show the integrability conditions are equivalent to the Gauss and Codazzi
  equations, $R_{ijkl}=\Pi_{jk}\Pi_{il}-\Pi_{ik}\Pi_{jl}$ and
  $(\nabla_i\Pi)_{jk}=(\nabla_j\Pi)_{ik}$; (3) conversely, given
  $g_{ij}$ and symmetric $\Pi_{ij}$ satisfying Gauss--Codazzi, show a local
  Riemannian immersion with second fundamental form $\Pi_{ij}$ exists; (4) for a
  metric of constant curvature $R^{-2}>0$, show there is a local Riemannian
  immersion into $\mathbb R^{n+1}$ landing in a sphere of radius $R$ (with unit
  normal $N=\pm R^{-1}F$).

  \emph{Indices.} The curvature equation in (2) is written here in the convention
  of \cref{thm:pet-ch3-tangential-curvature-equation}, with which it agrees: at a
  flat ambient that theorem reads $R^H(X,Y,Z,W)=\Pi(X,W)\Pi(Y,Z)-\Pi(X,Z)\Pi(Y,W)$,
  i.e. $R_{ijkl}=\Pi_{il}\Pi_{jk}-\Pi_{ik}\Pi_{jl}$. The printed source displays
  it as $R_{iklj}=\Pi_{ij}\Pi_{kl}-\Pi_{ik}\Pi_{jl}$, whose right-hand side is not
  antisymmetric in any pair of slots and so cannot be a curvature tensor as
  written; its first term matches, its second does not.
\end{remark}

\begin{remark}[Hyperbolic hypersurface realization]
  \label{rem:pet-ch3-ex-22}
  \lean{PetersenLib.exercise3_4_22}
  \source{petersen:page-0127}
  \uses{rem:pet-ch3-ex-21}
  \leanok
  \dcref{ch3:3.4.22}
  % Lean: the Gauss--Codazzi identities and converse constant-curvature implication are formalized; local existence is deferred.
  Repeat \cref{rem:pet-ch3-ex-21} for an immersion $F:M^n\hookrightarrow\mathbb
  R^{n,1}$ with normal $N$ satisfying $|N|^2=-1$, obtaining a local
  characterization of the hyperbolic spaces $H^n(R)$ as the local model of
  constant curvature $-R^{-2}$ (with unit normal $N=\pm R^{-1}F$). Only the sign
  $|N|^2=-1$ changes: the Gauss equation of (2) becomes
  $R_{ijkl}=-(\Pi_{jk}\Pi_{il}-\Pi_{ik}\Pi_{jl})$, while the Codazzi equation
  $(\nabla_i\Pi)_{jk}=(\nabla_j\Pi)_{ik}$ is unchanged.
\end{remark}

\begin{remark}[Kulkarni--Nomizu product]
  \label{rem:pet-ch3-ex-23}
  \lean{PetersenLib.exercise3_4_23}
  \source{petersen:page-0127}
  \uses{prop:pet-ch3-curvature-symmetries, def:pet-ch3-constant-curvature}
  \leanok
  \dcref{ch3:3.4.23}
  For symmetric $(0,2)$-tensors $h,k$, define
  $h\owedge k(v_1,v_2,v_3,v_4)=\tfrac12(h(v_1,v_4)k(v_2,v_3)+h(v_2,v_3)k(v_1,v_4))
  -\tfrac12(h(v_1,v_3)k(v_2,v_4)+h(v_2,v_4)k(v_1,v_3))$. (1) Show $h\owedge
  k=k\owedge h$. (2) Show $h\owedge h=0$ if $h$ has rank $1$. (3) Show that if
  $n>2$, $k$ is nondegenerate, and $h\owedge k=0$, then $h=0$. (4) Show $h\owedge
  k$ satisfies properties (1)--(3) of \cref{prop:pet-ch3-curvature-symmetries}.
  (5) Show $\nabla_X(h\owedge k)=(\nabla_Xh)\owedge k+h\owedge(\nabla_Xk)$. (6)
  Show $(M,g)$ has constant curvature $c$ iff the $(0,4)$-curvature tensor
  satisfies $R=c\cdot(g\owedge g)$.
\end{remark}

\begin{remark}[Schouten tensor]
  \label{rem:pet-ch3-ex-24}
  \lean{PetersenLib.exercise3_4_24}
  \source{petersen:page-0127,petersen:page-0128}
  \uses{def:pet-ch3-ricci-curvature, def:pet-ch3-scalar-curvature, rem:pet-ch3-ex-23}
  \leanok
  \dcref{ch3:3.4.24}
  Define $P=\dfrac2{n-2}\operatorname{Ric}-\dfrac{\operatorname{scal}}{(n-1)(n-2)}g$
  for $n>2$. (1) Show $P\equiv0$ implies $\operatorname{Ric}=0$. (2) Show the
  decomposition
  $P=\dfrac{\operatorname{scal}}{n(n-1)}g+\dfrac2{n-2}(\operatorname{Ric}-\tfrac{\operatorname{scal}}ng)$
  is orthogonal. (3) For $n=2$, show $R=\tfrac{\operatorname{scal}}2g\owedge g$.
  (4) For $n=3$, show
  $R=\tfrac{\operatorname{scal}}6g\owedge g+2(\operatorname{Ric}-\tfrac{\operatorname{scal}}3g)\owedge
  g=P\owedge g$. (5) For $n>2$, show $(M,g)$ has constant curvature iff
  $R=P\owedge g$ and $\operatorname{Ric}=\tfrac{\operatorname{scal}}ng$. (6) Show
  $\operatorname{Ric}(X,Y)=\sum_i(P\owedge g)(X,E_i,E_i,Y)$ for any orthonormal
  frame $E_i$.
\end{remark}

\begin{remark}[Weyl tensor]
  \label{rem:pet-ch3-ex-25}
  \lean{PetersenLib.exercise3_4_25}
  \source{petersen:page-0128}
  \uses{rem:pet-ch3-ex-24}
  \leanok
  \dcref{ch3:3.4.25}
  Define $W$ implicitly by
  $R=\dfrac{\operatorname{scal}}{n(n-1)}g\owedge g+\dfrac2{n-2}(\operatorname{Ric}-\tfrac{\operatorname{scal}}ng)\owedge
  g+W=P\owedge g+W$ (with $P$ as in \cref{rem:pet-ch3-ex-24}). (1) Show $n=3$
  implies $W=0$. (2) Show $\sum_iW(X,E_i,E_i,Y)=0$ for any orthonormal frame
  $E_i$. (3) Show the decomposition $R=P\owedge g+W$ is orthogonal.
\end{remark}

\begin{remark}[Divergence of the Schouten and Weyl tensors]
  \label{rem:pet-ch3-ex-26}
  \lean{PetersenLib.exercise3_4_26}
  \source{petersen:page-0129}
  \uses{rem:pet-ch3-ex-24, notation:pet-ch3-divergence-adjoint}
  \leanok
  \dcref{ch3:3.4.26}
  % Lean: the Schouten-divergence identity is formalized; the Weyl-divergence identity is deferred.
  Show $\nabla^*P=-\dfrac1{n-1}d\operatorname{scal}$ and
  $\nabla^*W(Z,X,Y)=\dfrac{n-3}2\big((\nabla_XP)(Y,Z)-(\nabla_YP)(X,Z)\big)$.
\end{remark}

\begin{remark}[Structure constants]
  \label{rem:pet-ch3-ex-27}
  \lean{PetersenLib.exercise3_4_27}
  \source{petersen:page-0129}
  \uses{def:pet-ch3-curvature-tensor}
  \leanok
  \dcref{ch3:3.4.27}
  For an orthonormal frame $E_1,\dots,E_n$ define structure constants $c^k_{ij}$
  by $[E_i,E_j]=c^k_{ij}E_k$ (functions on $M$), and $\Gamma^k_{ij},R^l_{ijk}$ by
  $\nabla_{E_i}E_j=\Gamma^k_{ij}E_k$, $R(E_i,E_j)E_k=R^l_{ijk}E_l$. Compute
  $\Gamma,R$ in terms of the $c^k_{ij}$; note that on a Lie group with a
  left-invariant metric the $c^k_{ij}$ may be taken constant.
\end{remark}

\begin{remark}[Cartan formalism]
  \label{rem:pet-ch3-ex-28}
  \lean{PetersenLib.exercise3_4_28}
  \source{petersen:page-0129,petersen:page-0130}
  \uses{def:pet-ch3-curvature-tensor}
  \leanok
  \dcref{ch3:3.4.28}
  % Lean: parts (1)--(2) are formalized; the converse in (1) and surface formula (3) are deferred.
  For a frame $E_1,\dots,E_n$, write the connection as $\nabla E_i=\omega^j_iE_j$
  (connection $1$-forms $\omega^j_i$), i.e.\ $\nabla_vE_i=\omega^j_i(v)E_j$; for
  an orthonormal frame with dual coframe $\omega^i$: (1) show $\omega^i_j=-\omega^j_i$
  and $d\omega^i=\omega^j\wedge\omega^i_j$, and that these two equations
  conversely determine the connection forms from the orthonormal frame. (2)
  Viewing $\omega^j_i$ as an $\mathfrak{so}(n)$-valued $1$-form, define the
  curvature forms $\Omega^j_i$ ($\mathfrak{so}(n)$-valued $2$-forms) by
  $R(X,Y)E_i=\Omega^j_i(X,Y)E_j$; show
  $d\omega^j_i=\omega^k_i\wedge\omega^j_k+\Omega^j_i$. (3) For surfaces, with
  orthonormal frame $E_1,E_2$ and coframe $\omega^1,\omega^2$, show
  $d\omega^1=\omega^2\wedge\omega^1_2$, $d\omega^2=-\omega^1\wedge\omega^1_2$,
  $d\omega^1_2=\Omega^1_2=\sec\cdot d\mathrm{vol}$.
\end{remark}

\begin{remark}[Polarization of algebraic curvature tensors]
  \label{rem:pet-ch3-ex-29}
  \lean{PetersenLib.exercise3_4_29}
  \source{petersen:page-0130,petersen:page-0131}
  \uses{lem:pet-ch3-bivector-jacobi-identity, prop:pet-ch3-curvature-symmetries}
  \leanok
  \dcref{ch3:3.4.29}
  Assume $R(X,Y,Z,W)$ is an algebraic curvature tensor, i.e.\ satisfies
  properties (1)--(3) of \cref{prop:pet-ch3-curvature-symmetries}. (1) Show
  \begin{align*}
    6R(X,Y,V,W)&=\frac{\partial^2}{\partial s\,\partial t}R(X+sW,Y+tV,Y+tW,X+sW)\Big|_{s=t=0}\\
    &\quad-\frac{\partial^2}{\partial s\,\partial t}R(X+sV,Y+tW,Y+tW,X+sV)\Big|_{s=t=0}.
  \end{align*}
  (2) Show $6R(X,Y,V,W)$ equals the explicit sum-of-twelve-terms expansion
  \begin{align*}
    &R(X+W,Y+V,Y+V,X+W)-R(X,Y+V,Y+V,X)-R(W,Y+V,Y+V,W)\\
    &\quad-R(X+W,V,V,X+W)-R(X+W,Y,Y,X+W)+R(X,V,V,X)+R(W,V,V,W)\\
    &\quad+R(X,Y,Y,X)+R(W,Y,Y,W)-R(X+V,Y+W,Y+W,X+V)\\
    &\quad+R(X,Y+W,Y+W,X)+R(V,Y+W,Y+W,V)+R(X+V,Y,Y,X+V)+R(X+V,W,W,X+V)\\
    &\quad-R(X,Y,Y,X)-R(V,Y,Y,V)-R(X,W,W,X)-R(V,W,W,V)
  \end{align*}
  (four terms are redundant, obtainable from polarizing $R(X,Y,Y,X)$).
\end{remark}

\begin{remark}[Curvature-operator norm from sectional curvature]
  \label{rem:pet-ch3-ex-30}
  \lean{PetersenLib.exercise3_4_30}
  \source{petersen:page-0131}
  \uses{rem:pet-ch3-ex-29, def:pet-ch3-curvature-operator}
  \leanok
  \dcref{ch3:3.4.30}
  Using \cref{rem:pet-ch3-ex-29}, show the norm of the curvature operator on
  $\Lambda^2T_pM$ is bounded, $|\mathfrak R|_p\le c(n)|\sec|_p$, for a constant
  $c(n)$ depending only on dimension, where $|\sec|_p$ is the largest absolute
  value of a sectional curvature at $p$.
\end{remark}

\begin{remark}[Curvature of left-invariant metrics]
  \label{rem:pet-ch3-ex-31}
  \lean{PetersenLib.exercise3_4_31}
  \source{petersen:page-0131}
  \uses{def:pet-ch3-curvature-tensor}
  \leanok
  \dcref{ch3:3.4.31}
  Let $G$ be a Lie group with a left-invariant nondegenerate metric
  $(\cdot,\cdot)$ on $\mathfrak g$; for $X\in\mathfrak g$ let
  $\operatorname{ad}_X^*$ be the adjoint of $\operatorname{ad}_XY=[X,Y]$ with
  respect to $(\cdot,\cdot)$, so $(\operatorname{ad}_X^*Y,Z)=(Y,[X,Z])$. For
  left-invariant fields:
  (1) the Levi-Civita connection is
  $\nabla_XY=\tfrac12([X,Y]-\operatorname{ad}_X^*Y-\operatorname{ad}_Y^*X)\in\mathfrak g$;
  it is torsion-free and metric-compatible.
  (2) $R(X,Y,Z,W)=-(\nabla_YZ,\nabla_XW)+(\nabla_XZ,\nabla_YW)-(\nabla_{[X,Y]}Z,W)$.
  (3)
  $R(X,Y,Y,X)=\tfrac14|\operatorname{ad}_X^*Y+\operatorname{ad}_Y^*X|^2-(\operatorname{ad}_X^*X,\operatorname{ad}_Y^*Y)-\tfrac34|[X,Y]|^2-\tfrac12([[X,Y],Y],X)-\tfrac12([[Y,X],X],Y)$.
  We treat the positive-definite (Riemannian) case. (Petersen's printed part~(1)
  has $+\operatorname{ad}_X^*Y$; with the adjoint convention above the correct
  sign is $-\operatorname{ad}_X^*Y$, as the biinvariant specialisation
  $\nabla_XY=\tfrac12[X,Y]$ confirms.)
\end{remark}

\begin{remark}[Curvature of biinvariant metrics]
  \label{rem:pet-ch3-ex-32}
  \lean{PetersenLib.exercise3_4_32}
  \source{petersen:page-0131,petersen:page-0132}
  \uses{rem:pet-ch3-ex-31, def:pet-ch3-curvature-operator}
  \dcref{ch3:3.4.32}
  Let $G$ be a Lie group with a biinvariant (possibly indefinite, nondegenerate)
  metric $(\cdot,\cdot)$ on $\mathfrak g$. Using left-invariant fields, show: (1)
  $\nabla_XY=\tfrac12[X,Y]$; (2) $R(X,Y)Z=\tfrac14[Z,[X,Y]]$; (3)
  $R(X,Y,Z,W)=-\tfrac14([X,Y],[Z,W])$, so sectional curvatures are nonnegative
  when $(\cdot,\cdot)$ is positive definite; (4) the curvature operator is
  nonnegative when $(\cdot,\cdot)$ is positive definite, by showing
  $g\big(\mathfrak R(\sum_{i=1}^kX_i\wedge Y_i),\sum_{i=1}^kX_i\wedge
  Y_i\big)=\tfrac14\big|\sum_{i=1}^k[X_i,Y_i]\big|^2$; (5) for
  $(\cdot,\cdot)$ positive
  definite, $\operatorname{Ric}(X,X)=0$ iff $X$ commutes with every left-invariant
  vector field, so $G$ has positive Ricci curvature if its center is discrete.
\end{remark}

\begin{remark}[Killing-form metric]
  \label{rem:pet-ch3-ex-33}
  \lean{PetersenLib.exercise3_4_33}
  \source{petersen:page-0132}
  \uses{rem:pet-ch3-ex-32, def:pet-ch3-ricci-curvature}
  \leanok
  \dcref{ch3:3.4.33}
  For a Lie group with nondegenerate Killing form $B$, using $-B$ as the
  left-invariant metric: (1) show this metric is biinvariant; (2) show
  $\operatorname{Ric}=-\tfrac14B$.
\end{remark}

\begin{remark}[Cartan formalism for Killing-form metrics]
  \label{rem:pet-ch3-ex-34}
  \lean{PetersenLib.exercise3_4_34}
  \source{petersen:page-0132}
  \uses{rem:pet-ch3-ex-33}
  \leanok
  \dcref{ch3:3.4.34}
  Using the Cartan formalism of \cref{rem:pet-ch3-ex-28} in the setting of
  \cref{rem:pet-ch3-ex-33}, compute all quantities in terms of the structure
  constants of $\mathfrak g$: with a biinvariant metric, in an orthonormal basis
  they satisfy $c^k_{ij}=-c^k_{ji}=c^i_{jk}$, the first equality being
  skew-symmetry of the bracket, the second biinvariance of the metric.
\end{remark}
